% Correspondencia editorial: gestion/mapa_transportes_auxiliares.md.
% Realizaciones auxiliares: no sustituyen el eje de compatibilidad y restitución.
\section{Transporte compensado, memoria geométrica y contraste energético}
\label{md:aux:transporte}

\subsection{Origen operatorio y dominio de la realización}
\label{md:aux:origen}

La construcción conserva APP--TRIT--TPK y el estado enriquecido. APP
reúne las hojas aritméticas con sus residuos y cocientes; TRIT conserva
el régimen y la orientación; TPK compone selección, transporte y
actualización de memoria. Las representaciones matriciales utilizadas
a continuación son realizaciones posteriores. La holonomía nonádica
$\Gamma_9$ y su memoria mantienen sus operadores propios.

Los antecedentes empleados son las representaciones cuadráticas
elíptica, hiperbólica y parabólica del TRIT, el diferencial de
transportes invertibles, el funcional energético
$E=\langle Df,WDf\rangle$ con $W>0$, el retorno con memoria
traslacional, la elipse de acción con producto de semiejes $2\hbar$
y su realización LC, la polarización de vacancias y el acoplamiento
reversible de dos modos. El retorno traslacional y el conmutador
que sigue son objetos distintos. El lazo elíptico--parabólico y
su balance de área orientada se utilizarán después como colección
de contraste.

La composición compensada y el protocolo de contraste reúnen estas
realizaciones. La teoría algebraica de trazas de conmutadores en
$\operatorname{SL}(2)$ constituye un antecedente matemático
establecido.\footnote{Goldman, \emph{Trace Coordinates on
Fricke spaces}, \href{https://arxiv.org/abs/0901.1404}{arXiv:0901.1404}.}
El objeto de este apéndice es su composición concreta con los
lectores declarados. Los parámetros $a,b$ son parámetros de las
realizaciones, no constantes físicas ajustadas.

El alcance termina en esta realización local. La identificación de
rutas TPK que admitan un dispositivo material con esos parámetros
requiere componer sus lectores efectivos. Una matriz invertible
arbitraria no acredita por sí sola una ruta TPK, y la identificación
del defecto de conmutación con torsión de Cartan requiere su
conexión y su soldadura.

\subsection{Compensación de operaciones y conservación de área}
\label{md:aux:conmutador}

Las matrices de régimen son
\begin{equation}
 J_{\mathrm e}=\begin{pmatrix}0&1\\-1&0\end{pmatrix},\qquad
 J_{\mathrm h}=\begin{pmatrix}0&1\\1&0\end{pmatrix}.
 \label{md:aux:regimenes}
\end{equation}
Satisfacen $J_{\mathrm e}^2=-I$, $J_{\mathrm h}^2=I$ y
$[J_{\mathrm e},J_{\mathrm h}]=2\operatorname{diag}(1,-1)$.
Se definen
\begin{equation}
 R(a)=e^{aJ_{\mathrm e}},\qquad B(b)=e^{bJ_{\mathrm h}},\qquad
 C_{\mathrm{aux}}(a,b)=R(a)B(b)R(a)^{-1}B(b)^{-1}.
 \label{md:aux:def-compensado}
\end{equation}
La acción sobre columnas se lee de derecha a izquierda. El protocolo
contiene cada transformación y su inversa, pero la sucesión no es
el retroceso completo del recorrido. Un retroceso completo tendría,
por ejemplo, la forma $RBB^{-1}R^{-1}=I$. El símbolo
$C_{\mathrm{aux}}$ distingue este conmutador elíptico--hiperbólico
del contraángulo $C^*$ y del lazo de formas conjugadas
$H_{\mathrm{lazo}}$ estudiado en el cuerpo principal.

\begin{proposition}
\label{md:aux:prop-traza}
Para $a,b\in\mathbb R$,
\begin{equation}
 \det C_{\mathrm{aux}}=1,\qquad
 \operatorname{tr}C_{\mathrm{aux}}=2+4\sin^2a\sinh^2b.
 \label{md:aux:traza-compensado}
\end{equation}
\end{proposition}
\begin{proof}
Póngase $c_a=\cos a$, $s_a=\sin a$, $u_b=\cosh b$ y $v_b=\sinh b$.
Entonces
\[
 R=\begin{pmatrix}c_a&s_a\\-s_a&c_a\end{pmatrix},\qquad
 B=\begin{pmatrix}u_b&v_b\\v_b&u_b\end{pmatrix},
\]
con $c_a^2+s_a^2=1$ y $u_b^2-v_b^2=1$. Todas las matrices tienen
determinante uno. La multiplicación da
\begin{equation}
 \begin{aligned}
 (C_{\mathrm{aux}})_{11}&=u_b^2+2c_as_au_bv_b-(c_a^2-s_a^2)v_b^2,\\
 (C_{\mathrm{aux}})_{12}&=-2s_a^2u_bv_b-2c_as_av_b^2,\\
 (C_{\mathrm{aux}})_{21}&=-2s_a^2u_bv_b+2c_as_av_b^2,\\
 (C_{\mathrm{aux}})_{22}&=u_b^2-2c_as_au_bv_b-(c_a^2-s_a^2)v_b^2.
 \end{aligned}
 \label{md:aux:entradas-compensado}
\end{equation}
La suma diagonal es
$2u_b^2-2(c_a^2-s_a^2)v_b^2=2+4s_a^2v_b^2$, como se afirma.
\end{proof}

Cuando $s_av_b\ne0$, la traza supera dos y el conmutador es
hiperbólico, con autovalores positivos recíprocos
\begin{equation}
 e^{\lambda},\ e^{-\lambda},\qquad
 \lambda=2\operatorname{arsinh}|\sin a\sinh b|>0.
 \label{md:aux:espectro-compensado}
\end{equation}
Se obtiene de $\cosh\lambda=\operatorname{tr}(C_{\mathrm{aux}})/2$
y $\cosh(2\operatorname{arsinh}x)=1+2x^2$.
Si $s_av_b=0$, entonces $b=0$ o $R=\pm I$, y
$C_{\mathrm{aux}}=I$.
La expansión local es
\begin{equation}
 C_{\mathrm{aux}}=I+2ab\operatorname{diag}(1,-1)
 +O(|a|^2|b|+|a||b|^2).
 \label{md:aux:expansion-compensado}
\end{equation}
El residuo comienza por el producto de las dos operaciones. Aplicar
cualquiera de ellas seguida de su propia inversa produce la identidad.

Una matriz real bidimensional de determinante uno conserva la forma
de área. Por ello, $C_{\mathrm{aux}}$ transforma la elipse de
acción en otra elipse de área $h$. Si $\Sigma>0$ es una matriz
de covarianza,
\begin{equation}
 \Sigma'=C_{\mathrm{aux}}\Sigma C_{\mathrm{aux}}^{\mathsf T},
 \qquad \det\Sigma'=\det\Sigma.
 \label{md:aux:covarianza-area}
\end{equation}
Pueden cambiar la distribución de semiejes y su orientación.
Los autovalores $e^{\pm\lambda}$ describen direcciones propias e
iteraciones; los valores singulares determinan otra lectura y, en
general, no coinciden con ellos. La variación de semiejes euclídeos
en una sola aplicación requiere la métrica correspondiente.

La conservación del área de acción es compatible con una
acumulación de deformación orientada. El observable geométrico
$\lambda$ depende de $a,b$ y no constituye una constante universal
adicional. La conservación del determinante o del área, por sí
sola, no establece conservación de energía física.

\subsection{Lector positivo del residuo e inversión del recorrido}
\label{md:aux:residuo}

En una fibra de retorno, con campo $f$ y peso hermítico $W>0$
declarados, el lector energético toma la forma
\begin{equation}
 E_C(f)=\langle(I-C_{\mathrm{aux}})f,
                  W(I-C_{\mathrm{aux}})f\rangle.
 \label{md:aux:energia-residuo}
\end{equation}
Para $s_av_b\ne0$, el conmutador carece de autovalor uno y
$E_C(f)>0$ para todo $f\ne0$. El retroceso completo produce residuo
cero. La normalización dimensional pertenece al funcional de
partida: sus coordenadas numéricas requieren su realización antes
de interpretarse como julios, calor, masa o corriente de conducción.

La inversión completa del recorrido invierte $C_{\mathrm{aux}}$.
La traza y $\lambda$ identifican el operador y su inverso. La
restitución de dirección requiere una lectura adicional: el vector
residual, las componentes de transporte o una sonda con resolución
de fase. Esta comparación limita de manera concreta la suficiencia
de una lectura exclusivamente escalar.

Un ejemplo racional, independiente de datos físicos, es
\begin{equation}
 R=\begin{pmatrix}3/5&4/5\\-4/5&3/5\end{pmatrix},\qquad
 B=\begin{pmatrix}5/4&3/4\\3/4&5/4\end{pmatrix}.
 \label{md:aux:ejemplo-racional}
\end{equation}
En él se obtiene $\operatorname{tr}C_{\mathrm{aux}}=86/25$,
$\det C_{\mathrm{aux}}=1$ y autovalores
$(43\pm6\sqrt{34})/25$. Los controles racionales del suplemento
incluyen este ejemplo y $169$ pares, con casos nulos, retroceso
completo y conservación del determinante de covarianza. Son
comprobaciones finitas de las identidades demostradas.

\subsection{Acoplamiento reversible y lectura reducida de entrelazamiento}
\label{md:aux:acoplamiento}

Para $H\in\operatorname{Sp}(2,\mathbb R)$, la realización de dos
modos utiliza el acoplamiento
\begin{equation}
 U_H=\frac19\begin{pmatrix}
 8I+H&\sqrt8(H-I)\\
 \sqrt8(H-I)&I+8H
 \end{pmatrix}.
 \label{md:aux:acoplamiento-uh}
\end{equation}
Con dos gaussianas puras idénticas preparadas desde la elipse de
acción y con transporte conjunto de la coordenatización inicial, el lector
reducido satisface
\begin{equation}
 \nu(H)^2=1+\frac8{81}
    \bigl\{\operatorname{tr}(HH^{\mathsf T})-2\bigr\},
 \qquad\operatorname{Tr}\rho_{\mathrm{vis}}^2=\nu(H)^{-1}.
 \label{md:aux:mezcla-generica}
\end{equation}
La preparación y los lectores se mantienen fijados. La
transformación conjunta es reversible y conserva la pureza global.

\begin{corollary}
\label{md:aux:cor-mezcla}
Para $C_{\mathrm{aux}}(a,b)$,
\begin{equation}
 \operatorname{tr}(C_{\mathrm{aux}}C_{\mathrm{aux}}^{\mathsf T})-2
 =16\sin^2a\sinh^2b\cosh^2b,
 \label{md:aux:traza-cuadratica}
\end{equation}
y
\begin{equation}
 \boxed{\nu_C^2=1+\frac{32}{81}\sin^2a\sinh^2(2b).}
 \label{md:aux:mezcla-compensado}
\end{equation}
\end{corollary}
\begin{proof}
Las cuatro entradas de \eqref{md:aux:entradas-compensado} se
escriben como $d+k$, $l-m$, $l+m$, $d-k$, donde
\[
 d=1+2s_a^2v_b^2,\quad k=2c_as_au_bv_b,\quad
 l=-2s_a^2u_bv_b,\quad m=2c_as_av_b^2.
\]
La suma de sus cuadrados es
$2(d^2+k^2+l^2+m^2)=2+16s_a^2u_b^2v_b^2$.
La sustitución en \eqref{md:aux:mezcla-generica} y
$\sinh(2b)=2\sinh b\cosh b$ prueban el resultado.
\end{proof}
El ejemplo \eqref{md:aux:ejemplo-racional} produce
$\nu_C^2=17/9$ y pureza $3/\sqrt{17}$.

La cancelación de $\hbar$ y de la deformación inicial utiliza
la preparación declarada y el transporte conjunto de coordenatización.
Un protocolo con $s_av_b\ne0$ produce mezcla reducida; la
inversión completa del acoplamiento puede deshacerla. Mezcla
reducida, energía almacenada y calor conservan sus funciones
como lecturas distintas.

\subsubsection{Contraste firmado en una colección elíptico--parabólica}
\label{md:aux:familia-hn}

La colección construida de lazos es
\begin{equation}
 H_n=\begin{pmatrix}1+n&n^2\\n&n^2-n+1\end{pmatrix},
 \qquad n\in\mathbb Z.
 \label{md:aux:hn}
\end{equation}
Su balance alternante de memoria depende de $n^2$, mientras el
cociente de determinantes de covarianza es
\begin{equation}
 \mathcal R_n=\frac{\det V_{\mathrm{vis},n}}{\det V_0}
 =1+\frac8{81}n^2(2n^2-2n+5).
 \label{md:aux:covarianza-hn}
\end{equation}
Al restar las dos orientaciones se cancelan los términos pares:
\begin{equation}
 \boxed{\mathcal R_{-n}-\mathcal R_n=\frac{32n^3}{81}.}
 \label{md:aux:diferencia-direccional}
\end{equation}
Para $n=1$ y $n=-1$, las respuestas son $121/81$ y $153/81$,
con purezas $9/11$ y $3/\sqrt{17}$, respectivamente.
Esta colección satisface $H_{-n}\ne H_n^{-1}$ en general;
tampoco es la iteración $H_1^n$. El cambio del parámetro
orientado debe distinguirse del retroceso completo de las
operaciones. Su lectura alternante de área conserva una
identificación distinta de la del funcional de Barbero.

El criterio de mezcla emplea el conmutador $[D,J]$, con
$D=\sqrt8(H-I)/9$. Si $D$ conmuta con la estructura compleja
compatible con el estado inicial, puede transportar memoria
sin mezclar ese producto gaussiano. Si no conmuta, el lector
determina la mezcla. La ausencia de mezcla para ese estado
es una propiedad distinta de la ausencia de disipación material.

La resta \eqref{md:aux:diferencia-direccional} es un corolario
de la colección y de su lector, no una constante universal.
La fórmula de $C_{\mathrm{aux}}$ compone el mismo lector con
otra colección admitida por el teorema. Un contraste material
debe determinar transportes y acoplamientos antes de medir;
la programación arbitraria de esas matrices acredita su
emulación y conserva un alcance diferente.

\subsection{Segundo momento de la corriente de vacancias}
\label{md:aux:vacancias}

El lector de capacidad fija
\[
 \delta=\log_{10}(10/9),\qquad z_n\in\{0,1\},\qquad
 P_N=\sum_{n<N}(z_n-\delta),\qquad |P_b-P_a|<1.
\]
En intervalos iguales $t_0$, la corriente centrada es
$I_n=A_I(z_n-\delta)$, con $A_I=u_Q/t_0$.
El símbolo $A_I$ distingue esta amplitud del ángulo principal.
Para $L=b-a$,
\begin{equation}
 \frac1{A_I^2}\sum_{n=a}^{b-1}I_n^2
 =L\delta(1-\delta)+(1-2\delta)(P_b-P_a).
 \label{md:aux:segundo-momento}
\end{equation}
La prueba consiste en desarrollar $(z_n-\delta)^2$, utilizar
$z_n^2=z_n$ y sumar. El desequilibrio firmado permanece
acotado, mientras
$I_{\mathrm{rms}}/|A_I|\longrightarrow\sqrt{\delta(1-\delta)}$.
La amplitud $A_I$ es el salto entre los dos niveles de corriente.
El contraste conserva la realización temporal y debe contabilizar
cualquier acoplamiento que filtre la sucesión.

Un balance firmado nulo puede coexistir con una lectura
cuadrática positiva. Si se añade una resistencia ideal $R_E>0$
y se mantiene esa corriente por intervalos, el modelo posterior
de contraste da
\begin{equation}
 Q_{R_E}=R_Et_0\sum_n I_n^2.
 \label{md:aux:contraste-resistivo}
\end{equation}
La resistencia pertenece a esta realización posterior. Su
función consiste en contrastar una posible reducción de calor
con un balance explícito.

\subsection{Condiciones de un contraste material}
\label{md:aux:contraste-material}

El protocolo separa dos cuestiones. Para la memoria de orden
se comparan recorridos compensados, intercambios de orden y
retrocesos completos mediante una respuesta remanente vectorial
o resuelta en fase. Una intensidad o una carga neta aisladas
pueden perder la diferencia buscada. Para la recuperación de
energía se miden entrada, energía útil o devuelta, variación
de almacenamiento y calor, incluidos preparación y restauración
del dispositivo. El retorno de fase conserva separada la
variación de energía almacenada.

En un contraste escalar de un puerto, la relación
$v_2(t)=v_1(T-t)$ conserva la distribución de amplitudes y el
espectro de potencia. En régimen periódico, un mismo sistema
lineal estacionario predice igual potencia media. Una ley
instantánea $i=g(v)$ también conserva la integral de $vi$.
Una diferencia reproducible examinaría dependencia histórica,
pero su atribución requiere comparar histéresis, calentamiento
y modelos dinámicos alternativos. En varios puertos deben
conservarse además los espectros cruzados.

Una reducción de disipación se contrasta a igual transferencia
útil, fidelidad y ritmo. La atenuación debe discriminar energía,
amplitud de señal, población de portadores y carga conservada
antes de atribuirle un origen gravitatorio. La distinción entre
almacenamiento, recuperación y pérdidas cuenta con precedentes
experimentales en recuperación electrocalórica.\footnote{Defay
et al., \emph{Enhanced electrocaloric efficiency via energy
recovery} (2018),
\href{https://doi.org/10.1038/s41467-018-04027-9}{doi:10.1038/s41467-018-04027-9}.
Este antecedente no constituye evidencia experimental del
transporte HMT aquí compuesto.}

El resultado del apéndice es una identidad general dentro
de la colección matricial especificada, un lector positivo del
residuo, un control de inversión y una ley exacta del segundo
momento de la corriente. Su alcance distingue esos resultados
de una nueva constante universal, un prototipo sin calentamiento
o una demostración de decaimiento electrónico gravitatorio.
La vinculación prospectiva con un dispositivo debe fijarse
antes de obtener los datos de contraste.

\section{Acción, memoria reducida y trabajo recuperable}
\label{md:aux:trabajo}

\subsection{Composición de la acción con una realización energética}
\label{md:aux:accion}

La composición recibe las salidas de acción y los transportes
de memoria de APP--TRIT--TPK, conservando el estado enriquecido
como antecedente. Su dominio es la realización gaussiana y
energética declarada. La pasividad y la ergotropía son conceptos
establecidos; se emplean para estudiar el acoplamiento específico
y los lectores ya construidos.

La ecuación de acción mantiene explícitos sus coeficientes:
\begin{equation}
 H_5=\frac12\sqrt{\frac{1000\alpha}{\varphi}}-\alpha
 +\frac9{16}\alpha^2-\frac59\alpha^3
 +\frac7{48}\alpha^4-\frac1{54}\alpha^5,
 \label{md:aux:h5}
\end{equation}
\begin{equation}
 D_A=e^{-100\pi\alpha/9},\qquad
 \mathcal C_\pi=169\alpha^6+\frac{D_A\alpha^7}{1-D_A\alpha},
 \label{md:aux:correccion-accion}
\end{equation}
\begin{equation}
 \begin{aligned}
 \hbar_{\mathrm{pre}}&=s_0H_5,\\
 \hbar_{\mathrm{ret}}&=s_0\left(H_5-\frac{90}{\pi}\mathcal C_\pi\right),
 \qquad h_s=2\pi\hbar_s,\qquad s_0=10^{-34}\mathcal U_S.
 \end{aligned}
 \label{md:aux:secciones-accion}
\end{equation}
La década y la coordenatización dimensional conservan antecedentes
distintos. La fórmula mantiene las coordenadas $\pi,\varphi$,
la orientación y la escala del lector, además de $\alpha$.
La igualdad dimensional conserva su unidad de acción.

Fijada una sección positiva, se escribe $\hbar=\hbar_s$.
La elipse tiene semiejes $\sqrt{2\hbar}e^{\zeta/2}$ y
$\sqrt{2\hbar}e^{-\zeta/2}$, donde
$\zeta=\operatorname{artanh}(C^*/A)$. Su área orientada cumple
\begin{equation}
 \mathcal A(\theta)=\hbar\theta,\qquad
 \mathcal A(2\pi)=h.
 \label{md:aux:area-fase}
\end{equation}
En la realización LC, con reloj $\omega$ e impedancia de
referencia declarados,
\begin{equation}
 \mathcal H_\zeta(Q,P)
 =\frac\omega2\left(e^{-\zeta}Q^2+e^\zeta P^2\right).
 \label{md:aux:hamiltoniano-lc}
\end{equation}
Sobre la órbita se obtiene $\mathcal H_\zeta=\hbar\omega=hf$,
con $f=\omega/(2\pi)$. La convención hamiltoniana da
$\dot\theta=-\omega$ y, por ello,
$d\mathcal A/dt=-\hbar\omega$. La energía positiva coincide
con la magnitud de esa velocidad areal en la realización y
sección fija consideradas. El área acumulada y la acción
lagrangiana de una trayectoria arbitraria son objetos distintos.

La ecuación interna determina la escala de acción por unidad
de fase; el reloj fija su ritmo. La deformación de la elipse
y el orden del transporte conservan información adicional a
$hf$. La preparación gaussiana empleada después es distinta
de un punto de la órbita clásica: su energía inicial es
$hf/2$, en lugar de $hf$.

\subsection{Preparación y lectores del resultado energético}
\label{md:aux:hipotesis-gaussianas}

Se preparan dos modos gaussianos puros, centrados e idénticos,
con covarianza
\begin{equation}
 V_0=\frac\hbar2MM^{\mathsf T},\qquad
 M=\operatorname{diag}(e^{\zeta/2},e^{-\zeta/2}).
 \label{md:aux:preparacion}
\end{equation}
Actúa $U_H$ de \eqref{md:aux:acoplamiento-uh}, transportado a
la coordenatización inicial mediante $M\oplus M$. El estado conjunto
permanece puro. Los lectores energéticos finales tienen una
misma frecuencia $\omega$, conservan la métrica inicial
$G_\zeta=M^{-\mathsf T}M^{-1}$ y no incluyen energía de
interacción residual entre los modos. Las operaciones de
extracción comienzan y terminan con esos mismos lectores.
La accesibilidad física de las operaciones admitidas es una
hipótesis de la cota de trabajo, distinta de la construcción
de un dispositivo que las realice.

Se definen
\begin{equation}
 \begin{gathered}
 \tau=\operatorname{tr}(HH^{\mathsf T})-2\geq0,\\
 W_v=\frac{8I+HH^{\mathsf T}}9,\qquad
 W_m=\frac{I+8HH^{\mathsf T}}9.
 \end{gathered}
 \label{md:aux:covarianzas-reducidas}
\end{equation}
Las covarianzas son $V_j=(\hbar/2)MW_jM^{\mathsf T}$.
La pureza visible se denota por $\mu=\operatorname{Tr}\rho_v^2$
y su autovalor simpléctico normalizado por $\nu=\mu^{-1}$.
En este apéndice, $\nu$ tiene esta función; la frecuencia
ordinaria se escribe $f$ y la angular $\omega$.

\subsection{Identidad entre energía visible y pureza}
\label{md:aux:energia-pureza}

\begin{proposition}
\label{md:aux:prop-energia-pureza}
Con la preparación y los lectores anteriores, la energía
visible por encima del estado inicial satisface
\begin{equation}
 \frac{\Delta E_v}{hf}=\frac\tau{36},\qquad
 \nu^2=1+\frac{8\tau}{81},\qquad
 \boxed{\mu^{-2}-1=\frac{32}{9}\frac{\Delta E_v}{hf}.}
 \label{md:aux:identidad-energia-pureza}
\end{equation}
\end{proposition}
\begin{proof}
El estado es centrado y el lector cuadrático permanece fijo.
Por tanto,
\[
 E_v=\frac\omega2\operatorname{tr}(G_\zeta V_v)
 =\frac{\hbar\omega}{4}\operatorname{tr}W_v.
\]
Como $\operatorname{tr}W_v=2+\tau/9$, al restar
$\hbar\omega/2$ se obtiene la primera igualdad.
La covarianza reducida satisface
$\det W_v=1+8\tau/81$ y $\mu=(\det W_v)^{-1/2}$.
La eliminación de $\tau$ prueba la identidad encuadrada.
La demostración no requiere simetría de $H$ ni coincidencia
entre sus autovalores y sus valores singulares.
\end{proof}

La ley elimina los parámetros particulares del recorrido.
El coeficiente $32/9$ procede del reparto del acoplamiento
y la escala $h$ se conserva hasta la normalización.
Para comprobar la especificidad de ese reparto, considérese
$W_p=pI+(1-p)HH^{\mathsf T}$, con $0<p<1$. Se obtiene
\begin{equation}
 \mu^{-2}-1=4p\frac{\Delta E}{hf}.
 \label{md:aux:reparto-general}
\end{equation}
La elección declarada es $p=8/9$. Cuando $\Delta E>0$,
un contraste puede estimar
\begin{equation}
 p_{\mathrm{obs}}=
 \frac{\mu^{-2}-1}{4\Delta E/(hf)}
 \label{md:aux:reparto-observado}
\end{equation}
y compararlo con $8/9$ sin ajustar $p$ a esos mismos datos.
Programar un mezclador con ese reparto verifica una
emulación; la atribución del reparto a un mecanismo natural
conserva su prueba independiente.

La identidad utiliza gaussianas centradas y la preparación
especificada. Un desplazamiento coherente añade energía sin
cambiar la pureza. El ruido térmico inicial y los cambios de
métrica alteran la relación. Estas variaciones delimitan el
enunciado y proporcionan falsadores de una extensión a estados
arbitrarios.

La combinación convexa $W_v$ es una combinación de matrices
de covarianza producida al reducir un estado gaussiano conjunto.
Una mezcla probabilística de dos estados puros constituye
otro objeto, generalmente no gaussiano; su pureza no viene
dada por el determinante de esa covarianza.

\subsection{Trabajo recuperable bajo operaciones locales}
\label{md:aux:ergotropia}

La ergotropía es el máximo descenso de energía mediante
operaciones unitarias cíclicas con lector energético fijo.
La definición no incluye automáticamente el coste neto de
construcción, preparación o control del aparato.

El espectro reducido tiene poblaciones
\begin{equation}
 \lambda_j=(1-r)r^j,\qquad r=\frac{\nu-1}{\nu+1}.
 \label{md:aux:espectro-gaussiano}
\end{equation}
Están ordenadas decrecientemente con la energía del oscilador.
El estado pasivo tiene así energía $E_{v,\mathrm{pas}}=hf\nu/2$.
Lo alcanza una operación gaussiana:
\begin{equation}
 S=\sqrt\nu\,W_v^{-1/2},\qquad
 \det S=1,\qquad SW_vS^{\mathsf T}=\nu I.
 \label{md:aux:pasivacion}
\end{equation}
El orden de las poblaciones demuestra también la minimalidad
entre todas las unitarias, además de las gaussianas.

Por consiguiente,
\begin{equation}
 \boxed{
 \frac{\mathcal W_v}{hf}
 =\frac{(\nu-1)(9\nu-7)}{32},\qquad
 \frac{E_{v,\mathrm{pas}}-hf/2}{hf}=\frac{\nu-1}{2}.}
 \label{md:aux:trabajo-visible}
\end{equation}
Las dos cantidades suman
$\Delta E_v/(hf)=9(\nu^2-1)/32$. El término pasivo es
energía inaccesible bajo las operaciones locales declaradas;
su presencia no equivale a un flujo de calor realizado.

\subsubsection{Acceso conjunto a las correlaciones}
\label{md:aux:acceso-conjunto}

El segundo modo cumple $\det W_m=\det W_v=\nu^2$ y
$\Delta E_m=8\Delta E_v$. Por ello,
\begin{equation}
 \Delta E_{\mathrm{tot}}=9\Delta E_v=hf\frac\tau4.
 \label{md:aux:energia-conjunta}
\end{equation}
La inversión conjunta del acoplamiento restituye el producto
inicial. El estado total es puro y ese producto es el estado
fundamental del lector conjunto; el acceso conjunto permite
recuperar idealmente $\Delta E_{\mathrm{tot}}$.
Las operaciones locales independientes dejan dos estados
pasivos, cada uno de energía $hf\nu/2$. Se deduce
\begin{equation}
 \boxed{\mathcal W_{\mathrm{conjunta}}-\mathcal W_{\mathrm{local}}
 =hf(\nu-1)=hf(\mu^{-1}-1).}
 \label{md:aux:ventaja-conjunta}
\end{equation}

La identidad cuantifica el trabajo adicional accesible
mediante las correlaciones. La preparación inicial requirió
energía y el balance de un ciclo completo incluye ese aporte.
La relación compara lecturas energéticas y operaciones sobre
el estado, sin identificar información y energía como una
misma magnitud física.

En este enunciado, ``local'' significa unitarias de producto
$U_v\otimes U_m$; quedan fuera de esa clase las mediciones,
la comunicación clásica y la realimentación. El incremento
procede de la restricción operativa sobre el estado
correlacionado y de sus espectros marginales, no de una
energía de interacción residual. El reparto $1:8$ concierne
a las energías por encima del vacío, no a las energías
totales ni a las ergotropías.

\subsection{Casos exactos y orientación del transporte}
\label{md:aux:casos-energia}

Para $H_n$ de \eqref{md:aux:hn},
\begin{equation}
 \tau_n=n^2(2n^2-2n+5).
 \label{md:aux:tau-hn}
\end{equation}
En $n=1$ se obtiene
\begin{equation}
 \begin{gathered}
 \nu=11/9,\qquad\mu=9/11,\qquad
 \Delta E_v=\frac5{36}hf,\\
 \mathcal W_v=\frac1{36}hf,\qquad
 E_{v,\mathrm{pas}}-hf/2=\frac19hf.
 \end{gathered}
 \label{md:aux:caso-positivo}
\end{equation}
El conjunto tiene $\Delta E_{\mathrm{tot}}=5hf/4$, trabajo
local $37hf/36$ y ventaja conjunta $2hf/9$.
El caso $n=-1$ comparte los autovalores de $H_1$,
pero tiene $\tau=9$ y
\begin{equation}
 \nu=\frac{\sqrt{17}}3,\qquad
 \Delta E_v=\frac{hf}{4},\qquad
 \mathcal W_v=\frac{9-2\sqrt{17}}{12}hf.
 \label{md:aux:caso-negativo}
\end{equation}
La energía visible añadida difiere en razón $9/5$ a
frecuencia y sección de acción iguales. El cambio $n\mapsto-n$
no es el retroceso completo: $H_{-n}\ne H_n^{-1}$ en general.
La inversión matricial verdadera conserva $\tau$ en
dimensión dos. La asimetría paramétrica de esta colección
conserva así su distinción respecto de una reversión completa.

Para $C_{\mathrm{aux}}$,
$\tau_C=4\sin^2a\sinh^2(2b)$. La composición energética da
\begin{equation}
 \boxed{\Delta E_v(C_{\mathrm{aux}})
 =\frac{hf}{9}\sin^2a\sinh^2(2b),}
 \label{md:aux:energia-compensado}
\end{equation}
\begin{equation}
 \mathcal W_{\mathrm{conjunta}}-\mathcal W_{\mathrm{local}}
 =hf\left(\sqrt{1+\frac{32}{81}\sin^2a\sinh^2(2b)}-1\right).
 \label{md:aux:trabajo-compensado}
\end{equation}
La conexión con Planck permanece explícita: en la sección
previa, $hf=2\pi s_0H_5(\alpha,\varphi)f$; en la de retorno,
$H_5$ se sustituye por $\eta_{\mathrm{ret}}$. La composición
mantiene fijada la sección utilizada; no supone una
variación experimental de $\alpha$ ni una alternancia libre
entre secciones dimensionales.

\subsection{Boltzmann y temperatura equivalente del estado pasivo}
\label{md:aux:temperatura}

Para $\nu>1$, el estado pasivo tiene temperatura equivalente
\begin{equation}
 k_BT_{\mathrm{pas}}
 =\frac{hf}{\log[(\nu+1)/(\nu-1)]}.
 \label{md:aux:temperatura-pasiva}
\end{equation}
Cuando se utiliza el reloj del transductor,
$hf=k_B\Theta_{\mathrm{clk}}\log3$, resulta
\begin{equation}
 \frac{T_{\mathrm{pas}}}{\Theta_{\mathrm{clk}}}
 =\frac{\log3}{\log[(\nu+1)/(\nu-1)]}.
 \label{md:aux:razon-temperaturas}
\end{equation}
Para $n=1$ se recupera $\log3/\log10$. La entropía es
la del espectro gaussiano, con ocupación $(\nu-1)/2$.
La composición energética conserva esta lectura térmica
previamente construida. La temperatura equivalente del
estado pasivo no implica que el estado reducido previo
estuviese en equilibrio ni que se haya transferido calor.

\subsection{Especificidad y condiciones de contraste}
\label{md:aux:contraste-trabajo}

La relación general entre coherencia, correlaciones y
extracción de trabajo pertenece a la literatura existente.
La composición examinada especifica una ley conjunta:
coeficiente $32/9$, reparto energético $1:8$, colección
orientada $H_n$, escala de acción y diferencia entre
acceso local y conjunto.

Se distinguen dos niveles experimentales. La emulación
realiza $U_H$ y comprueba sus identidades; verifica esa
realización del modelo. La predicción física exige
determinar prospectivamente, desde la construcción, el
dispositivo, los grados de libertad, la frecuencia, la
preparación y el acoplamiento que realizan sus objetos,
sin ajustar $8/9$ a los datos del contraste. Se miden
por separado energía, covarianza, pureza, costes de
control y energía de restauración.

Un contraste independiente obtiene temperatura y pureza
por lectores independientes de la ley que pretende
comprobar. El trabajo recuperado se compara con un
balance que incluya preparación, controles y restauración
de la memoria.

Las comparaciones externas incluyen estudios de
ergotropía gaussiana y pasivación,\footnote{Kua, Serafini
y Genoni, \emph{Daemonic ergotropy of Gaussian quantum
states and the role of measurement-induced purification
via general-dyne detection},
\href{https://arxiv.org/abs/2506.22288}{arXiv:2506.22288}.
Su papel es la comparación con fórmulas generales de
energía, pureza y pasividad; no valida el acoplamiento HMT
ni sus constantes.}
la extracción de ergotropía coherente con costes
energéticos explícitos,\footnote{\emph{Experimental
Extraction of Coherent Ergotropy and Its Energetic Cost
in a Superconducting Qubit},
\href{https://arxiv.org/abs/2506.16881}{arXiv:2506.16881}.}
y el control coherente entre ciclos de máquinas
cuánticas.\footnote{Hou et al., \emph{Combining energy
efficiency and quantum advantage in cyclic machines},
\href{https://doi.org/10.1038/s41467-025-60179-5}{doi:10.1038/s41467-025-60179-5}.
El experimento citado no mide el transporte nonádico
compuesto en este apéndice.}

El resultado es una composición algebraica con
estructura de partida explícita. Sus pruebas fijan
identidades y restricciones operativas. Una ventaja
experimental observada, una nueva constante universal,
una violación de la termodinámica o un mecanismo material
de transporte sin calor requerirían sus respectivas
demostraciones y contrastes; no son conclusiones de las
identidades auxiliares aquí establecidas.
